% This supplementary chapter preserves the full account of Pardon's derived
% regularity argument. It is not part of the default fourteen-talk route.

\chapter{Pardon's derived regularity argument}
\label[chapter]{chap:Pardon_Derived_Regularity}

Chapters 12--14 establish representability for elliptic moduli problems and
apply it to pseudo-holomorphic maps. Their route is local and analytic. Near a
chosen solution, one passes to Sobolev completions, adds a finite-dimensional
obstruction space, constructs a Kuranishi chart, and proves that this chart
represents an open restriction of the intrinsic derived solution stack.

Pardon's argument reorganizes the same ingredients. It begins with ordinary
Fredholm theory only on the locus where the linearized equation is
surjective. On this regular locus, classical analysis produces an ordinary
smooth moduli manifold. The new question is then:

\begin{quote}
  Why does the same smooth manifold classify regular families over derived
  test objects?
\end{quote}

The answer is a theorem about smooth sections over compact sources. Their
stack on derived manifolds is the sheaf left Kan extension of the
corresponding stack on ordinary manifolds. Since sheaf left Kan extension
preserves pullbacks along submersive maps, the ordinary regular moduli
manifold also represents the regular derived moduli stack. Finally, every
nonregular solution becomes regular after enlarging the parameter space by
finitely many obstruction directions. The original moduli problem is
recovered as the derived zero-parameter fiber of this enlarged regular
problem.

Thus the analytic and categorical parts of the proof are sharply separated.
All elliptic estimates concern ordinary smooth parameter spaces. The passage
to derived parameters is formal once the section-stack theorem is available.

There are two sources, with different status. Pardon's published 2024
proceedings paper states the pseudo-holomorphic Derived Regularity Theorem and
gives a proof sketch
\cite[Sections~4--5]{Pardon_Representability_Elliptic_Fredholm}. His July
2026 manuscript gives a much fuller general elliptic treatment, but explicitly
labels itself unfinished work in progress and retains unresolved internal
references
\cite[P.5 and N.6]{Pardon_Logarithmic_Derived_Moduli}. We will use the
published statement as the endpoint, explain the proof architecture in
detail, and mark precisely where the longer manuscript remains provisional.

\section{From a singular zero to the regularity theorem}
\label[section]{sec:Pardon_Singular_Zero_Regularity_Theorem}

\subsection{A regular enlargement of \texorpdfstring{$x^2=0$}{x squared equals zero}}
\label[subsection]{sec:Pardon_X_Squared_Enlargement}

The whole argument is already visible in a one-dimensional example. Consider
the smooth function
\[
  f\colon\R\longrightarrow\R,
  \qquad
  f(x)=x^2.
\]
Its ordinary zero locus is the point $\{0\}$, but the derivative at the zero
vanishes. The derived zero locus remembers this failure of transversality. Its
tangent complex at the origin is
\[
  [\R\xrightarrow{0}\R],
\]
with the first copy in deformation degree and the second in obstruction
degree.

Add a parameter $t$ and define
\begin{equation}
  \label{eq:Pardon_X_Squared_Thickening}
  \widetilde f\colon\R^2\longrightarrow\R,
  \qquad
  \widetilde f(x,t)=x^2+t.
\end{equation}
The derivative of $\widetilde f$ at $(0,0)$ is projection onto the second
coordinate, hence is surjective. Its zero locus
\[
  V=\widetilde f^{-1}(0)
  =\{(x,t):t=-x^2\}
\]
is therefore a smooth one-dimensional manifold. Restrict the parameter
coordinate to obtain
\[
  \kappa\colon V\longrightarrow\R,
  \qquad
  \kappa(x,-x^2)=-x^2.
\]
The original derived zero locus is recovered as
\begin{equation}
  \label{eq:Pardon_X_Squared_Derived_Fiber}
  \R\mathbin{\times^R}_{\R}\{0\}
  \iso
  V\mathbin{\times^R}_{\R}\{0\},
\end{equation}
where the map on the left is $f$ and the map on the right is $\kappa$. The
left side presents the singular equation directly. The right side presents it
as the derived zero-parameter fiber of the regular enlarged equation
\Cref{eq:Pardon_X_Squared_Thickening}.

The general finite-dimensional pattern is identical. Suppose
$f\colon\R^n\to\R^m$ satisfies $f(x)=0$. Choose a linear map
\[
  \alpha\colon\R^k\longrightarrow\R^m
\]
whose composite onto $\operatorname{coker}(D_xf)$ is surjective, and put
\[
  \widetilde f(y,t)=f(y)+\alpha(t).
\]
Then $D_{(x,0)}\widetilde f$ is surjective. After shrinking,
$\widetilde f^{-1}(0)$ is a manifold, and the derived zero locus of $f$ is
its derived fiber over the zero value of the added parameter.

This construction performs two tasks. The enlarged zero locus is accessible
to the ordinary implicit function theorem, while the final derived fiber
retains the obstruction data of the original equation. Pardon's elliptic
argument is the family version of this picture.

\subsection{The intrinsic elliptic solution stack}
\label[subsection]{sec:Pardon_Intrinsic_Elliptic_Solution_Stack}

Let
\[
  W\longrightarrow C\longrightarrow B
\]
be a family of section problems. We first take $B$ to be a smooth manifold,
$C\to B$ to be a proper submersion, and the differential equation to be
elliptic. The pseudo-holomorphic case from
\Cref{chap:Pseudo_Holomorphic_Curves_Derived_Moduli} is the running example.

Pardon defines the solution stack diagrammatically. For a test object $T$, a
$T$-point is a map $T\to B$ together with a section
\[
  u\colon C\times_BT\longrightarrow W
\]
which solves the equation in the vertical $C/B$-direction. This definition
makes sense for smooth, derived smooth, and suitable topological test objects.
In the derived case, solving the equation is additional equality data in a
mapping anima, not a property that can be checked only at ordinary points.

Locally near a chosen section, the section-stack charts of Chapter~9 and a
trivialization of the output bundle put the equation into the form
\begin{equation}
  \label{eq:Pardon_Equation_Morphism}
  \Phi\colon
  \operatorname{Sect}_{C/B}(W)
  \longrightarrow
  \operatorname{Sect}(C_0,H_0),
\end{equation}
where $C\cong B\times C_0$ by Ehresmann's theorem and $C_0$ is compact. The
right-hand side is the stack of sections of a fixed vector bundle
$H_0\to C_0$. The solution stack is the fiber over the zero section:
\begin{equation}
  \label{eq:Pardon_Solution_Stack_Pullback}
  \begin{tikzcd}[column sep=large]
    \operatorname{Hol}_B(C,W)
      \rar \dar \drar[pullback]
      & \operatorname{Sect}_{C/B}(W) \dar["\Phi"] \\
    * \rar["0"']
      & \operatorname{Sect}(C_0,H_0).
  \end{tikzcd}
\end{equation}
The point of the local form is not to replace the intrinsic definition. It
exhibits that definition as a pullback of two section stacks to which the
categorical extension theorem will apply.

For pseudo-holomorphic sections, $H_0$ is locally the bundle of vertical
complex-antilinear one-forms and $\Phi(u)=\bar\partial_Ju$. The diagrammatic
definition does not use the problematic full first-jet target discussed in
\Cref{warn:Steffens_Five_Tuple_Fourteen}.

\subsection{Regularity in a family}
\label[subsection]{sec:Pardon_Regularity_In_Family}

At a solution $(b,u)$, the total linearization of
\Cref{eq:Pardon_Equation_Morphism} has the form
\begin{equation}
  \label{eq:Pardon_Total_Linearization}
  D_{(b,u)}\Phi\colon
  \Gamma(C_b,u^*T_{W/C})\oplus T_bB
  \longrightarrow
  \Gamma(C_b,u^*H).
\end{equation}
The first summand varies the section and the second varies the parameter in
the family. In the fixed-domain pseudo-holomorphic case, the first summand is
the Cauchy--Riemann operator $D_u$ from
\Cref{prop:Linearized_Cauchy_Riemann_Operator_Fourteen}.

\begin{definition}[Regular locus]
  \label[definition]{def:Pardon_Regular_Locus}
  The \emph{regular locus}
  $\operatorname{Hol}_B(C,W)^{\mathrm{reg}}$ consists of the solutions
  $(b,u)$ at which \Cref{eq:Pardon_Total_Linearization} is surjective.
\end{definition}

Surjectivity is open for Fredholm operators, so this is an open substack of
the solution stack. The basic analytic input is ordinary parametric
regularity.

\begin{theorem}[Ordinary regularity]
  \label[theorem]{thm:Pardon_Ordinary_Regularity}
  Let $B$ be a smooth manifold and let $W\to C\to B$ be a smooth family of
  elliptic section problems with $C\to B$ proper. Then the regular locus
  $\operatorname{Hol}_B(C,W)^{\mathrm{reg}}$ on smooth test manifolds is
  represented by a smooth manifold, locally over $B$. Its comparison with
  the corresponding topological regular locus is an isomorphism.
\end{theorem}

This is the first step of the proof of Pardon's published Theorem~5.1. It is
stated as P.5.1 and proved in N.5.2.1 of the longer manuscript
\cite[Theorem~5.1]{Pardon_Representability_Elliptic_Fredholm}
\cite[P.5.1 and N.5.2.1]{Pardon_Logarithmic_Derived_Moduli}. Its analytic
ingredients are Newton--Picard iteration, elliptic regularity, and smooth
dependence on ordinary parameters. They are the regular version of the
Fredholm reduction developed in Chapter~10.

The theorem is stronger than a pointwise implicit-function statement. It is
compatible with smooth base change in the target section stack. This
parametric form will show that $\Phi$ is a submersive map of smooth stacks on
the regular locus.

\subsection{The published theorem}
\label[subsection]{sec:Pardon_Published_Derived_Regularity}

We can now state the endpoint of the argument.

\begin{theorem}[Derived Regularity Theorem, published form]
  \label[theorem]{thm:Pardon_Published_Derived_Regularity}
  Let $W\to C\to B$ be a pseudo-holomorphic section problem over a derived
  smooth stack $B$. Suppose that these maps are submersions, that $C\to B$ is
  proper of relative dimension two, and that the relevant vertical tangent
  bundles carry the prescribed complex structures. Then
  \[
    \operatorname{Hol}_B(C,W)\longrightarrow B
  \]
  is representable. Moreover, the tautological comparison
  \[
    (\mathrm{Der}\to\mathrm{Top})_!
      \operatorname{Hol}_B(C,W)_{\mathrm{Der}}
    \longrightarrow
      \operatorname{Hol}_B(C,W)_{\mathrm{Top}}
  \]
  is an isomorphism.
\end{theorem}

This is Pardon's Theorem~5.1
\cite[Theorem~5.1]{Pardon_Representability_Elliptic_Fredholm}. The printed
statement says that $B\to C$ is proper of relative dimension two. Since the
composable maps are $W\to C\to B$, the intended proper morphism is
$C\to B$.

The theorem concerns smooth compact domains. It does not add nodal curves or
compactify the moduli stack. Its conclusion does not explicitly include
quasi-smoothness. In the two-term pseudo-holomorphic problem, the tangent
complex calculated in
\Cref{eq:Relative_Tangent_Complex_Final_Fourteen} gives quasi-smoothness
once the represented object's geometric tangent complex is identified with
that analytic complex.

The published article calls its argument a brief proof sketch. We will now
prove the smooth-base deduction from two named inputs: ordinary regularity and
the proper section-stack theorem. This is the conceptual core of Pardon's
proof and avoids doing analysis on derived parameter objects.

\section{Proper section stacks and smooth tests}
\label[section]{sec:Pardon_Proper_Section_Stacks}

\subsection{The universal extension problem}
\label[subsection]{sec:Pardon_Universal_Extension_Problem}

Let
\[
  i\colon\mathrm{Sm}\longrightarrow\mathrm{Der}
\]
be the inclusion of smooth manifolds into derived smooth manifolds. Restriction
along $i$ gives a functor
\[
  i^*\colon
  \operatorname{Shv}(\mathrm{Der})
  \longrightarrow
  \operatorname{Shv}(\mathrm{Sm}).
\]
It has a left adjoint
\[
  i_!\colon
  \operatorname{Shv}(\mathrm{Sm})
  \longrightarrow
  \operatorname{Shv}(\mathrm{Der}),
\]
called \emph{sheaf left Kan extension}. Thus, for a smooth stack $F$ and a
derived smooth stack $X$, there is a natural isomorphism of mapping animae
\[
  \Hom_{\operatorname{Shv}(\mathrm{Der})}(i_!F,X)
  \iso
  \Hom_{\operatorname{Shv}(\mathrm{Sm})}(F,i^*X).
\]

This adjunction characterizes $i_!F$ as the universal derived smooth stack
built from $F$ while respecting colimits and descent. It does not merely assign
the set of ordinary families to a derived test object. Its values are animae,
and the colimit construction retains all coherence forced by derived
presentations.

For any derived smooth stack $X$, the adjunction gives a comparison morphism
\begin{equation}
  \label{eq:Pardon_Left_Kan_Comparison_General}
  i_!i^*X\longrightarrow X.
\end{equation}
There is no reason for
\Cref{eq:Pardon_Left_Kan_Comparison_General} to be an isomorphism for an
arbitrary stack. The theorem below says that section stacks over proper
families have this special property.

\subsection{The section-stack extension theorem}
\label[subsection]{sec:Pardon_Section_Stack_Extension_Theorem}

\begin{proposition}[Proper section stacks]
  \label[proposition]{prop:Pardon_Proper_Section_Stack_Extension}
  Let $W\to C\to B$ be submersions of smooth manifolds and suppose that
  $C\to B$ is proper. Then the canonical comparison
  \begin{equation}
    \label{eq:Pardon_Proper_Section_Stack_Comparison}
    i_!\operatorname{Sect}_{C/B}(W)_{\mathrm{Sm}}
    \longrightarrow
    \operatorname{Sect}_{C/B}(W)_{\mathrm{Der}}
  \end{equation}
  is an isomorphism, in the relative sense over $B$.
\end{proposition}

For $B=*$, this says that the derived stack of sections over a compact
manifold is determined by its restriction to ordinary smooth test manifolds.
For general $B$, the precise formulation uses the site of smooth manifolds
submersive over $B$. This relative version is exactly what is needed in
\Cref{eq:Pardon_Solution_Stack_Pullback}.

This is Proposition~5.2 of the published paper
\cite[Proposition~5.2]{Pardon_Representability_Elliptic_Fredholm}. The longer
manuscript separates the fixed-source and relative forms as T.7.7.2 and
T.7.7.4
\cite[T.7.7.2 and T.7.7.4]{Pardon_Logarithmic_Derived_Moduli}.

Several qualifications are essential.

\begin{enumerate}
  \item The proposition is about the full stack of smooth sections. No
    differential equation has yet been imposed.
  \item It does not say that derived families are ordinary families. It says
    that their functor on derived tests is the universal coherent extension
    of the smooth functor.
  \item It is not a formal consequence of the construction of derived
    manifolds. A general derived smooth stack need not satisfy
    \Cref{eq:Pardon_Left_Kan_Comparison_General}.
  \item Properness and partitions of unity are the geometric reasons the
    statement holds.
\end{enumerate}

\subsection{Why properness and softness enter}
\label[subsection]{sec:Pardon_Properness_Softness}

We give the proof architecture for $B=*$. Let $C$ be compact. A derived
manifold is locally a finite limit of ordinary manifolds. More precisely, the
essential-image criterion used by Pardon says that a derived smooth stack is
left Kan extended from manifolds if it carries finite cosifted presentations
of derived manifolds to the corresponding relative limit diagrams.

Suppose that
\[
  Q\iso\lim_{\alpha}Q_\alpha
\]
is such a finite presentation by derived manifolds, locally refined by smooth
pieces. Pulling back the bundle $W\to C$ gives a diagram over
\[
  Q\times C\iso\lim_{\alpha}(Q_\alpha\times C).
\]
Before taking global sections, the universal property of derived manifolds
shows that the resulting sheaves of local sections form the required relative
limit diagram.

Global sections are obtained by pushing those local sheaves forward along
the projection
\[
  Q\times C\longrightarrow Q.
\]
This is the step at which compactness matters. Proper base change controls the
pushforward because $C$ is compact Hausdorff. After passing to an open
section-stack chart, one may suppose that $W\to C$ is a vector bundle. Its
sheaf of smooth sections is soft: partitions of unity extend and glue local
sections. Softness implies that the proper pushforward preserves the finite
colimit diagrams occurring in the essential-image criterion. Consequently,
the global section stack satisfies that criterion and is left Kan extended
from smooth manifolds.

The relative proof works in the same way over the site of manifolds
submersive over $B$. This discussion is a proof sketch, not a replacement for
the higher-topos details in Pardon's sources. Its purpose is to identify the
mathematical jobs of the hypotheses:

\begin{enumerate}
  \item The finite-limit description of derived manifolds provides the
    presentations on which the extension criterion is tested.
  \item Properness controls pushforward from the source family.
  \item Partitions of unity make the relevant sheaves soft.
  \item Softness makes the necessary colimit diagrams survive pushforward.
\end{enumerate}

The smooth setting is crucial. Pardon points out that the analogous statement
fails in analytic geometry under the same compactness hypothesis, precisely
because the required softness is absent
\cite[proof of Proposition~5.2]{Pardon_Representability_Elliptic_Fredholm}.

\subsection{Which pullbacks survive left Kan extension?}
\label[subsection]{sec:Pardon_Pullbacks_Left_Kan_Extension}

The solution stack is a pullback of section stacks, but
\Cref{prop:Pardon_Proper_Section_Stack_Extension} does not yet determine it.
A left adjoint preserves colimits, not arbitrary limits. In particular, sheaf
left Kan extension need not preserve the pullback
\Cref{eq:Pardon_Solution_Stack_Pullback}.

The missing input is the following categorical principle.

\begin{lemma}[Left Kan extension and submersive pullbacks]
  \label[lemma]{lem:Pardon_Left_Kan_Submersive_Pullbacks}
  Suppose a functor between smooth geometric sites preserves pullbacks along
  submersions. Then its sheaf left Kan extension preserves pullback squares in
  which the distinguished morphism is submersive.
\end{lemma}

This is the specialization of Pardon's Lemma~5.3 used in the proof
\cite[Lemma~5.3]{Pardon_Representability_Elliptic_Fredholm}. At the
representable level, the assertion is the assumed preservation of submersive
pullbacks. General sheaves are assembled from representables by colimits, and
colimits in an $\infty$-topos are stable under base change. The representable
base-change isomorphisms therefore extend to the relevant pullback squares of
presheaves. The topological-site hypotheses and sheafification give the sheaf
version.

For the inclusion $i\colon\mathrm{Sm}\to\mathrm{Der}$, the representable
input is familiar from Chapters~3--4: a pullback along a submersion is
transverse, so its ordinary smooth pullback already computes the derived
pullback. The lemma promotes this fact from representable objects to sheaves.

The direction of the condition matters. In the regular-locus square below,
the zero section need not be submersive. It is the equation morphism on the
other side of the square which will be submersive.

\section{The regular-locus pullback}
\label[section]{sec:Pardon_Regular_Locus_Pullback}

\subsection{The equation morphism is submersive on the regular locus}
\label[subsection]{sec:Pardon_Equation_Submersive_Regular_Locus}

Restrict the source of \Cref{eq:Pardon_Equation_Morphism} to the open locus
where its linearization is surjective:
\[
  \Phi^{\mathrm{reg}}\colon
  \operatorname{Sect}_{C/B}(W)^{\mathrm{reg}}
  \longrightarrow
  \operatorname{Sect}(C_0,H_0).
\]
To show that this is submersive as a morphism of smooth stacks, use the
testing criterion recalled in
\Cref{prop:Smooth_To_Derived_Submersion_Criterion_Twelve}. Let $Z$ be a manifold and
let
\[
  \psi\colon Z\longrightarrow\operatorname{Sect}(C_0,H_0)
\]
be a smooth family of inhomogeneous terms. Form the smooth pullback
\begin{equation}
  \label{eq:Pardon_Test_Pullback}
  \mathcal P_Z
  =
  \operatorname{Sect}_{C/B}(W)^{\mathrm{reg}}
  \times_{\operatorname{Sect}(C_0,H_0)}Z.
\end{equation}
A point of this pullback is a triple $(b,u,z)$ satisfying
\[
  \Phi(b,u)=\psi(z).
\]
This is another family of elliptic equations on the compact source $C_0$.
Ordinary parametric regularity represents its regular locus by a smooth
manifold.

The derivative of the equation at a solution is
\begin{equation}
  \label{eq:Pardon_Base_Changed_Linearization}
  (\xi,\beta,v)
  \longmapsto
  D_{(b,u)}\Phi(\xi,\beta)-D_z\psi(v).
\end{equation}
Here $\xi$ varies the section, $\beta\in T_bB$ varies the original parameter,
and $v\in T_zZ$ varies the new parameter. The tangent space to the represented
solution manifold is the kernel of
\Cref{eq:Pardon_Base_Changed_Linearization}.

\begin{lemma}[Kernel projection]
  \label[lemma]{lem:Pardon_Kernel_Projection}
  If $D_{(b,u)}\Phi$ is surjective, then the projection
  \[
    \ker\bigl(D_{(b,u)}\Phi-D_z\psi\bigr)
    \longrightarrow T_zZ
  \]
  is surjective.
\end{lemma}

\begin{proof}
  Let $v\in T_zZ$. By surjectivity of $D_{(b,u)}\Phi$, choose
  $(\xi,\beta)$ such that
  \[
    D_{(b,u)}\Phi(\xi,\beta)=D_z\psi(v).
  \]
  Then $(\xi,\beta,v)$ lies in the displayed kernel and projects to $v$.
\end{proof}

The ordinary submersion theorem now shows that the represented pullback
\Cref{eq:Pardon_Test_Pullback} is submersive over $Z$. Since this holds after
every base change from a manifold, $\Phi^{\mathrm{reg}}$ is a submersive map of
smooth stacks.

This elementary calculation deserves emphasis. N.6.4.1 of the 2026
manuscript identifies the same submersive locus by invoking the comparison
between analytic and geometric tangent complexes, whose proof is still marked
by an unresolved reference. Here we use only the tangent-space conclusion of
the ordinary parametric implicit function theorem. Thus the calculation
reconstructs the required smooth submersiveness without assuming that missing
derived tangent comparison.

\subsection{The regular-locus square}
\label[subsection]{sec:Pardon_Regular_Locus_Square}

Restricting \Cref{eq:Pardon_Solution_Stack_Pullback} gives
\begin{equation}
  \label{eq:Pardon_Regular_Locus_Pullback_Square}
  \begin{tikzcd}[column sep=large]
    \operatorname{Hol}_B(C,W)^{\mathrm{reg}}
      \rar \dar \drar[pullback]
      & \operatorname{Sect}_{C/B}(W)^{\mathrm{reg}}
        \dar["\Phi^{\mathrm{reg}}"] \\
    * \rar["0"']
      & \operatorname{Sect}(C_0,H_0).
  \end{tikzcd}
\end{equation}
We know three facts about this square.

\begin{enumerate}
  \item The two section stacks are left Kan extended from smooth manifolds by
    \Cref{prop:Pardon_Proper_Section_Stack_Extension}. The same remains true
    after restriction to the open regular locus.
  \item The point $*$ is represented by a smooth manifold and is therefore
    left Kan extended from smooth manifolds.
  \item The right vertical morphism is submersive by the preceding
    subsection.
\end{enumerate}

Apply \Cref{lem:Pardon_Left_Kan_Submersive_Pullbacks} to
\Cref{eq:Pardon_Regular_Locus_Pullback_Square}. We obtain the central
comparison.

\begin{proposition}[Regular families over derived tests]
  \label[proposition]{prop:Pardon_Regular_Locus_Extension}
  The canonical morphism
  \begin{equation}
    \label{eq:Pardon_Regular_Locus_Comparison}
    i_!\operatorname{Hol}_B(C,W)^{\mathrm{reg}}_{\mathrm{Sm}}
    \longrightarrow
    \operatorname{Hol}_B(C,W)^{\mathrm{reg}}_{\mathrm{Der}}
  \end{equation}
  is an isomorphism.
\end{proposition}

\begin{proof}
  Apply $i_!$ to the smooth version of
  \Cref{eq:Pardon_Regular_Locus_Pullback_Square}. The section-stack theorem
  identifies the images of the upper-right and lower-right corners with the
  corresponding derived section stacks. The lower-left corner remains $*$.
  Since the right vertical morphism is submersive, the pullback lemma says that
  the image square is still a pullback. Its upper-left corner is therefore the
  derived regular solution stack, which proves
  \Cref{eq:Pardon_Regular_Locus_Comparison}.
\end{proof}

By \Cref{thm:Pardon_Ordinary_Regularity}, the smooth stack on the left of
\Cref{eq:Pardon_Regular_Locus_Comparison} is represented by an ordinary smooth
manifold $M^{\mathrm{reg}}$. Its sheaf left Kan extension is represented by the
same manifold, now regarded as a derived manifold. We have proved:

\begin{corollary}[Derived representability on the regular locus]
  \label[corollary]{cor:Pardon_Derived_Regular_Locus_Represented}
  The derived regular locus
  $\operatorname{Hol}_B(C,W)^{\mathrm{reg}}_{\mathrm{Der}}$ is represented by
  the ordinary smooth manifold which represents its restriction to smooth test
  manifolds.
\end{corollary}

This is the conceptual climax of Pardon's argument. The proof never asks what
a Sobolev completion parameterized by a derived manifold should mean. All
analysis was performed after base change by ordinary manifolds. Properness,
softness, and the preservation of submersive pullbacks then transported the
answer to derived test objects.

\subsection{What the regularity hypothesis did}
\label[subsection]{sec:Pardon_What_Regularity_Did}

It is useful to locate regularity precisely. The section-stack theorem holds
on the entire stack of sections. The obstruction appears when we impose the
equation by a pullback. Sheaf left Kan extension need not preserve this
pullback in general. Surjectivity of the linearization makes the equation
morphism submersive, and that is exactly the class of pullbacks which the
extension functor preserves.

Thus regularity is not used to define the derived solution stack. Nor is it
needed to give that stack a tangent complex. It is used to prove that a
particular ordinary smooth representative continues to represent the same
functor on derived tests.

This also explains why the regular locus alone is insufficient. A singular
solution is a perfectly valid ordinary point of the moduli functor, but it is
not contained in the locus where the equation morphism is submersive. The
next step changes the family, not the solution, so that the same point becomes
regular in an enlarged problem.

\section{Removing regularity by parameter thickening}
\label[section]{sec:Pardon_Parameter_Thickening}

\subsection{Choosing obstruction directions}
\label[subsection]{sec:Pardon_Choosing_Obstruction_Directions}

Fix an arbitrary solution $(b,u)$. The vertical linearization of the elliptic
equation is Fredholm, so its cokernel is finite-dimensional. Choose a
finite-dimensional vector space $E$ and a linear map
\begin{equation}
  \label{eq:Pardon_Obstruction_Directions}
  \alpha\colon
  E\longrightarrow\Gamma(C_b,u^*H)
\end{equation}
whose composite to the cokernel is surjective. One may take $E$ to be the
cokernel itself and choose representatives of a basis.

Using local trivializations and partitions of unity, extend these
representatives to smooth inhomogeneous terms for nearby sections and nearby
base points. In the pseudo-holomorphic case, these are finitely many smooth
complex-antilinear one-forms along the chosen map, extended to a neighborhood
of its graph.

Enlarge the base to
\[
  B^+=B\times E
\]
and define a new family of equations schematically by
\begin{equation}
  \label{eq:Pardon_Enlarged_Equation}
  \Phi^+(b',u',e)=\Phi(b',u')+\alpha_{(b',u')}(e).
\end{equation}
The construction is arranged so that $\alpha_{(b',u')}(0)=0$. Hence the
restriction of the enlarged family along
\[
  B=B\times\{0\}\longrightarrow B^+
\]
is the original family of equations.

At $(b,u,0)$ the new linearization is
\begin{equation}
  \label{eq:Pardon_Augmented_Linearization}
  D\Phi^+(\xi,\beta,e)
  =D_{(b,u)}\Phi(\xi,\beta)+\alpha(e).
\end{equation}

\begin{lemma}[Finite-dimensional regularization]
  \label[lemma]{lem:Pardon_Finite_Dimensional_Regularization}
  The map in \Cref{eq:Pardon_Augmented_Linearization} is surjective.
\end{lemma}

\begin{proof}
  Let $y$ lie in the target. By the choice of $\alpha$, choose $e\in E$ such
  that $y-\alpha(e)$ has trivial image in
  $\operatorname{coker}(D_{(b,u)}\Phi)$. Thus
  $y-\alpha(e)=D_{(b,u)}\Phi(\xi,\beta)$ for some $(\xi,\beta)$, and
  \Cref{eq:Pardon_Augmented_Linearization} sends $(\xi,\beta,e)$ to $y$.
\end{proof}

The chosen solution is therefore regular in the enlarged family. This is the
same Fredholm linear algebra as the obstruction-space construction of
Chapter~10. The difference is organizational: the obstruction coefficients
are now coordinates on the enlarged base.

\subsection{The original problem as a derived fiber}
\label[subsection]{sec:Pardon_Original_Problem_Derived_Fiber}

Let
\[
  \operatorname{Hol}_{B^+}(C^+,W^+)^{\mathrm{reg}}
\]
denote the regular locus of the enlarged problem. It contains the point
$(b,u,0)$ and, by
\Cref{cor:Pardon_Derived_Regular_Locus_Represented}, is represented near that
point by an ordinary smooth manifold $M^+_{\mathrm{reg}}$.

Pull back along the zero-parameter inclusion $B\to B^+$. Formation of the
intrinsic moduli stack is compatible with base change, so
\begin{equation}
  \label{eq:Pardon_Local_Derived_Chart}
  M^+_{\mathrm{reg}}
  \mathbin{\times^R}_{B^+}B
  \longrightarrow
  \operatorname{Hol}_B(C,W)
\end{equation}
represents an open neighborhood of $(b,u)$. The pullback must be formed in
derived manifolds. Although $M^+_{\mathrm{reg}}$ is smooth, its map to
$B^+$ need not be transverse to the zero-parameter inclusion.

The parallel with \Cref{eq:Pardon_X_Squared_Derived_Fiber} is exact:
\[
  \begin{array}{c@{\qquad}c}
    V\mathbin{\times^R}_{E}\{0\}
    &
    M^+_{\mathrm{reg}}
      \mathbin{\times^R}_{B^+}B.\\[4pt]
    \text{finite-dimensional model}
    &
    \text{elliptic moduli problem}
  \end{array}
\]
In both cases, a regular smooth manifold is first constructed in an enlarged
family. Its derived zero-parameter fiber is the singular object we wanted.

In the two-term situation, the local chart in
\Cref{eq:Pardon_Local_Derived_Chart} is quasi-smooth. Locally the inclusion
$B\to B\times E$ is the zero section of a finite-rank vector bundle, so the
relative tangent complex of the pullback is represented by a two-term complex
with obstruction term $E$.

\subsection{Covering the full moduli stack}
\label[subsection]{sec:Pardon_Covering_Full_Moduli_Stack}

The construction applies at every solution. For each $(b,u)$, choose finitely
many obstruction directions and an enlarged family in which that point is
regular. Regularity is open, so pulling the enlarged regular locus back to
$B$ gives an open substack of the original moduli stack containing $(b,u)$.
These open substacks cover the underlying topological solution stack.

Representability is local. Since every member of the cover is represented by
a finite derived pullback of manifolds, the full derived solution stack is
represented locally by derived manifolds. There is no separate compatibility
problem among the choices of obstruction directions. Each chart represents
an open restriction of the same intrinsic moduli functor, and uniqueness of
representing objects supplies the transition data.

The comparison with the topological moduli stack is proved in parallel.
Ordinary regularity identifies the topology on each enlarged regular locus,
and the zero-parameter fibers cover the original moduli stack. Since being an
isomorphism is local, the comparison is an isomorphism globally.

We can package the argument in a form which separates the completed deduction
from the status of its inputs.

\begin{theorem}[Derived representability from regularity]
  \label[theorem]{thm:Pardon_Representability_From_Regularity}
  Let $B$ be a smooth manifold and let $W\to C\to B$ be a family of two-term
  elliptic section problems with $C\to B$ proper. Assume the following.
  \begin{enumerate}
    \item Ordinary parametric regularity represents every regular locus and
      its smooth-to-topological comparison.
    \item The relevant section stacks satisfy the comparison in
      \Cref{prop:Pardon_Proper_Section_Stack_Extension}.
  \end{enumerate}
  Then the intrinsic derived solution stack is locally represented by
  quasi-smooth derived manifolds, and its derived-to-topological comparison is
  an isomorphism.
\end{theorem}

\begin{proof}
  By \Cref{prop:Pardon_Regular_Locus_Extension}, the regular locus of every
  enlarged problem is represented on derived test objects by the same smooth
  manifold that represents it on smooth tests. The construction of
  \Cref{sec:Pardon_Choosing_Obstruction_Directions} makes every chosen solution
  regular after a finite-dimensional enlargement of the base.
  \Cref{eq:Pardon_Local_Derived_Chart} then represents an open neighborhood of
  that solution by a quasi-smooth derived manifold. These neighborhoods cover,
  so locality proves representability. The same cover and ordinary
  smooth-to-topological comparison prove the last assertion.
\end{proof}

For the pseudo-holomorphic problem, this is the three-step architecture of
\Cref{thm:Pardon_Published_Derived_Regularity}: ordinary regularity, extension
of the regular locus to derived tests, and parameter thickening.

\section{Comparison, scope, and source status}
\label[section]{sec:Pardon_Comparison_Scope_Status}

\subsection{Comparison with Steffens}
\label[subsection]{sec:Pardon_Comparison_With_Steffens}

Pardon and Steffens begin with the same intrinsic goal and the same Fredholm
linear algebra. Their proofs differ in where they place the smooth-to-derived
bridge.

\begin{enumerate}
  \item \emph{Intrinsic object.} Steffens forms the solution stack from a
    differential moduli problem. Pardon defines the pseudo-holomorphic
    solution stack diagrammatically and then expresses it locally as a fiber
    of section stacks.
  \item \emph{Analytic input.} Steffens constructs a finite-dimensional
    reduction near each solution by a Sobolev argument. Pardon first applies
    ordinary parametric regularity only on the surjective locus.
  \item \emph{Obstruction directions.} Steffens adds a finite-dimensional
    obstruction space to the operator. Pardon regards its coefficients as new
    coordinates on an enlarged parameter space. The augmented linearizations
    are the same.
  \item \emph{Smooth-to-derived bridge.} Steffens proves that one augmented
    operator is submersive in smooth stacks and transports that fact to
    derived stacks. Pardon proves the structural
    \Cref{prop:Pardon_Proper_Section_Stack_Extension} for all proper section
    stacks, then applies it to a submersive pullback on the regular locus.
  \item \emph{Local chart.} Steffens obtains the derived zero locus of a
    finite-dimensional obstruction section. Pardon obtains the derived
    zero-parameter fiber of an ordinary regular moduli manifold. These are two
    descriptions of the same local Kuranishi pattern.
  \item \emph{Globalization.} Both arguments cover the intrinsic moduli stack
    by represented open substacks and invoke locality. Neither needs an
    explicit cocycle of independently chosen Kuranishi charts.
\end{enumerate}

The main conceptual contribution of Pardon's organization is now visible:

\begin{quote}
  Flexible mapping objects over compact smooth sources are already controlled
  by their values on ordinary smooth parameter spaces.
\end{quote}

This principle is more general than the pseudo-holomorphic application, but
it is also special to smooth geometry. Properness without partitions of unity
does not give the same conclusion in analytic geometry.

\subsection{Representability is not compactification}
\label[subsection]{sec:Pardon_Representability_Not_Compactification}

The properness used throughout the proof belongs to the family of domains
$C\to B$. It ensures compact fibers and makes section-stack pushforward
behave well. It does not imply that the solution stack is compact or proper.

For pseudo-holomorphic curves, sequences of solutions can bubble and their
domains can degenerate to nodal curves. Such limiting domains are not objects
of the stack of smooth compact Riemann surfaces used here. The theorem therefore
does not supply the compactification needed for Gromov--Witten theory.

Representability also does not construct a virtual fundamental class. Even
for a compact quasi-smooth derived moduli stack, one still needs orientation
data and a theory which converts the derived structure into a homology or
bordism class. Pardon's implicit-atlas theory addresses such enumerative
questions by a different construction. The derived regularity argument should
not be described as a reformulation of implicit atlases.

The logical ledger is:

\begin{enumerate}
  \item Elliptic regularity and proper domains give finite-dimensional local
    deformation theory.
  \item Derived representability packages that theory coherently.
  \item Compactification adds missing limiting objects.
  \item Orientations and a virtual-class formalism produce enumerative data.
\end{enumerate}

Only the first two items have been addressed in this chapter.

\section{Summary and supplementary materials}
\label[section]{sec:Pardon_Derived_Regularity_Summary_Supplementary_Materials}

\subsection{The live route}
\label[subsection]{sec:Pardon_Live_Route}

The 90-minute talk can be compressed into the following statements.

\begin{enumerate}
  \item The singular zero locus of $x^2$ is the derived zero-parameter fiber
    of the regular equation $x^2+t=0$.
  \item An elliptic solution stack is intrinsically a fiber of section
    stacks.
  \item Ordinary Fredholm analysis represents the locus on which the total
    linearization is surjective.
  \item Sheaf left Kan extension is the universal passage from smooth stacks
    to derived smooth stacks.
  \item Proper section stacks are left Kan extended from smooth manifolds.
  \item Properness, softness, and partitions of unity are the geometric inputs
    to this extension theorem.
  \item Left Kan extension does not preserve arbitrary pullbacks.
  \item It does preserve pullbacks along submersive morphisms.
  \item On the regular locus, the equation morphism is submersive.
  \item Consequently, the same ordinary manifold represents regular families
    on smooth and derived test objects.
  \item Every nonregular solution becomes regular after adding finitely many
    obstruction parameters.
  \item The original moduli stack is the derived zero-parameter fiber of the
    enlarged regular problem.
  \item This yields local derived representability without analysis on derived
    parameter spaces.
  \item The result concerns smooth compact domains, not nodal compactification
    or virtual fundamental classes.
\end{enumerate}

The next chapter crosses precisely this boundary. It asks what happens when
the domain itself degenerates and why logarithmic smooth geometry is proposed
as a setting for the corresponding gluing problem.

\subsection{Supplement: the general pullback lemma}
\label[subsection]{sec:Pardon_Supplement_General_Pullback_Lemma}

Pardon's Lemma~5.3 is more general than
\Cref{lem:Pardon_Left_Kan_Submersive_Pullbacks}. Let
$f\colon C\to D$ be a functor, and let $P$ and $Q$ be pullback-stable classes
of morphisms in $C$ and $D$. Suppose $f$ takes pullback squares along
$P$-morphisms to pullback squares along $Q$-morphisms. Then
\[
  f_!\colon\operatorname{PSh}(C)\longrightarrow\operatorname{PSh}(D)
\]
has the same property. If $C$ and $D$ are perfect topological
$\infty$-sites and $f$ is topological, the corresponding statement holds for
sheaf left Kan extension
\[
  f_!\colon\operatorname{Shv}(C)\longrightarrow\operatorname{Shv}(D).
\]

The proof uses the expression of presheaves as colimits of representables and
the universality of colimits in presheaf and sheaf $\infty$-topoi. The
application takes $C=\mathrm{Sm}$, $D=\mathrm{Der}$, and both $P$ and $Q$ to
be submersions. The inclusion preserves the corresponding pullbacks because
ordinary pullbacks along submersions are transverse.

\subsection{Supplement: the 2026 proof and its unresolved inputs}
\label[subsection]{sec:Pardon_Supplement_Source_Audit}

The July 2026 manuscript materially expands the published proof sketch.
P.5.1 states ordinary regularity, P.5.2 states derived regularity, and N.6.4.1
gives the detailed passage from one to the other. The proper section-stack
theorems are T.7.7.2 and T.7.7.4, while N.6.2.1 gives the parameter-thickening
cover
\cite[P.5.1--P.5.4, T.7.7.1--T.7.7.4, and N.6.2--N.6.4]
  {Pardon_Logarithmic_Derived_Moduli}.

The manuscript nevertheless remains unfinished. Three unresolved points are
relevant here.

\begin{enumerate}
  \item In P.5.3, the comparison between the geometric tangent complex of the
    represented moduli object and the analytic complex of linearized elliptic
    operators is followed by ``(??)'' rather than a completed reference.
  \item N.6.4.1 invokes this comparison to recognize a base-changed regular
    locus as the submersive locus over a smooth parameter manifold. The
    kernel-projection argument in
    \Cref{lem:Pardon_Kernel_Projection} supplies the smooth parametric
    conclusion needed for the core proof, but it does not establish the full
    derived tangent comparison claimed in P.5.3.
  \item N.6.3.1 reduces elliptic problems over derived bases to problems over
    smooth bases by invoking a derived Ehresmann theorem marked ``(??)''.
    Consequently, our detailed proved route works over a smooth manifold or,
    by descent, a smooth stack base. The further derived-base reduction is
    reported as part of Pardon's work-in-progress generalization.
\end{enumerate}

There are further placeholders inside the flexibility results used to prove
T.7.7.2 and T.7.7.4. This does not erase the conceptual value of the argument,
but it does determine how its status should be described. The 2024 article is
a published proof sketch of the pseudo-holomorphic theorem. The 2026 manuscript
is a substantially fuller, explicitly unfinished general treatment, not yet a
source in which every dependency has been closed.

\subsection{Related literature}
\label[subsection]{sec:Pardon_Derived_Regularity_Related_Literature}

The primary source for the live route is Pardon's
\emph{Representability in non-linear elliptic Fredholm analysis}. Section~4
defines pseudo-holomorphic moduli stacks diagrammatically. Section~5 states the
Derived Regularity Theorem, describes its three proof steps, proves the
section-stack theorem in sketch form, and states the categorical pullback lemma
\cite[Sections~4--5]{Pardon_Representability_Elliptic_Fredholm}.

The July 2026 manuscript
\emph{Logarithmic derived moduli theory of non-linear elliptic equations}
places the same argument in a much larger treatment. The relevant ordinary
smooth-domain material is in P.4--P.5, T.7.7, and N.5--N.6.4. Its logarithmic
extensions begin afterwards and belong to the next chapter
\cite[P.4--P.5, T.7.7, and N.5--N.6.4]
  {Pardon_Logarithmic_Derived_Moduli}.

For comparison, Steffens's Remark~3.4.7 isolates the smooth-to-derived problem
inside the Kuranishi-chart proof. Chapter~11 followed that route in detail
\cite[Remark~3.4.7]{Steffens_Representability_Elliptic_Moduli}. The two
arguments use the same Fredholm obstruction spaces but package the categorical
comparison differently.
